\chapter{Introduction}

Power-system small-signal stability is a theory of operators before it is a
theory of plots. A linearized power system is stable or unstable because the
spectrum of its differential-algebraic representation lies in particular
regions of the complex plane, because the corresponding eigenvectors project
onto physical states in particular ways, and because perturbations of
parameters move those eigenvalues through mechanisms that can be traced
mathematically. The familiar language of damping ratio, participation factors,
critical modes, and modal frequencies is a compact engineering vocabulary for
that operator theory \cite{kundur1994power,rogers2000oscillations,kundur2004definition}.

The same vocabulary becomes harder to use in converter-rich systems. A
synchronous-machine dominated model permits a useful separation: inertia is
stored in rotating masses, synchronizing stiffness comes from the network and
power-angle relation, and damping can often be represented by a fixed local
coefficient. Converter-interfaced resources alter this separation. A
grid-following converter may contribute phase-locked-loop states, current-loop
states, plant controllers, filters, and protection-related dynamics. A
grid-forming converter may contribute droop, virtual oscillator, virtual
synchronous-machine, voltage-loop, current-loop, and current-limiting dynamics.
The extended IEEE/CIGRE stability taxonomy explicitly recognizes that
power-electronic interfaces introduce faster and device-dependent dynamic
phenomena that do not fit cleanly into the classical electromechanical
classification \cite{hatziargyriou2021stability}.

This dissertation asks a sharper theoretical question:

\begin{quote}
\itshape
What is the reduced dynamic operator that preserves grid strength, fragility,
and actionable pole motion in a converter-rich power system?
\end{quote}

That question has a planning consequence. If a utility asks how much
converter-based generation can be hosted, where virtual inertia should be
placed, or which lines should be reinforced before installing batteries, the
answer cannot be read from voltage limits or graph connectivity alone. A
hosting decision is dynamic only if it states the damping, resolvent, or
protection margin that must survive the conversion. A line decision is robust
only if it separates useful stiffness increase from damping-ratio loss or
controller-resonant reinforcement risk.

The answer developed here is hierarchical, but the hierarchy is not the final
claim. Fixed scalar damping is the uniform-decay member of a broader sequence
of increasingly expressive operators. The exact decoupling condition for fixed
damping is not \(D=\gamma M\) alone. It is commutation of the
inertia-normalized stiffness and damping operators. A fixed but non-commuting
\(D\) keeps the problem quadratic, but it couples graph modes through the modal
damping matrix. A controller-condensed converter-rich model changes the problem
more fundamentally: the reduced correction depends on the spectral parameter
\(s\). The effective operator is then a rational self-energy \(\Pi(s)\),
obtained by a Schur complement of the eliminated controller states. The object
that turns this hierarchy into an action theory is the dynamic Green function
\(G_T(s)=T(s)^{-1}\), whose port projections measure how finite physical
actions move poles.

\section{Thesis Statement}

The thesis statement is:

\begin{quote}
Dynamic grid strength in converter-rich power systems is an operator-valued
property. Its static limit is encoded by graph stiffness, its oscillatory
behavior by the Green function of the QEP, its controller-mediated behavior by
rational self-energy \(\Pi(s)\), its finite-action behavior by low-rank
determinant identities, its fragility by structured pseudospectra, and its
nonlinear extension by memory kernels. The damping-channel symbol
\(\Sigma(s)\) is reserved for the velocity channel after
\(\Pi(s)=\Pi_K(s)+s\Pi_D(s)\) has been declared. Spectral graph methods remain
valuable, but their role changes from certificate to screen whenever damping
and stiffness cease to share an eigenbasis, whenever eliminated controllers
contribute frequency-dependent dynamics, or whenever an action claim requires
a certified finite perturbation rather than a local derivative.
\end{quote}

This statement does not claim novelty for Laplacian linearizations,
generalized graph modes, effective resistance, QEP or NEP sensitivity
formulas, Schur complements, or the existence of Braess-like phenomena in
oscillator and power networks. Those are established mathematical and
engineering ingredients. The novelty sought by this theoretical thesis is the
synthesis: a single operator hierarchy that explains when classical damping
intuition is mathematically justified, when it fails, and what reduced object
must replace it.

\thesisbox{thesisorange}{Originality boundary}{
Known mathematics is used as the foundation; it is not claimed as the
dissertation's invention. The doctoral contribution is the operator-valued
planning synthesis. The strongest original claims are the graph roughness of
\(D_i/M_i\) as a modal-coupling obstruction, the graph-time response atlas,
the spectral effective resistance law for finite topology updates, and the
controller-induced reinforcement-reversal certificate. Claims about universal
weak nodes, completed sheaf-grid theory, or empirical benchmark closure remain
validation targets rather than defended theorems.}

\thesisbox{thesisgreen}{Doctoral reading contract}{
Every planning statement in this dissertation must answer four questions:
which operator is being used, which port or action family is being tested,
which margin is being protected, and which validation gate turns the
calculation from a screen into evidence.}

\section{The Operator Hierarchy}

The hierarchy begins with the second-order swing equation
\begin{equation}
  M\ddot\theta+D\dot\theta+L\theta=p,
  \label{eq:intro_swing}
\end{equation}
where \(M\) is inertia, \(D\) is damping, \(L\) is synchronizing stiffness, and
\(\theta\) is the vector of bus or machine angles after removing the uniform
reference. In inertia-normalized coordinates, the stiffness becomes
\(\Lt=M^{-1/2}LM^{-1/2}\) and the damping becomes
\(\Dt=M^{-1/2}DM^{-1/2}\).

The first theoretical floor is fixed damping. In inertia-normalized
coordinates,
\begin{equation}
  \Lt=M^{-1/2}LM^{-1/2},\qquad
  \Dt=M^{-1/2}DM^{-1/2}.
  \label{eq:intro_normalized_operators}
\end{equation}
Exact modal decoupling occurs when
\begin{equation}
  [\Lt,\Dt]=\Lt\Dt-\Dt\Lt=0.
  \label{eq:intro_commutation_condition}
\end{equation}
For symmetric matrices, this is the condition for simultaneous orthogonal
diagonalization. The familiar \(D=\gamma M\) case gives
\(\Dt=\gamma I\), so it certainly satisfies
\eqref{eq:intro_commutation_condition}; it is the special case in which every
underdamped mode has the same real decay rate \(-\gamma/2\). Rayleigh damping
\(D=\alpha M+\beta L\) also commutes in normalized coordinates and therefore
decouples graph modes, but it does not produce a uniform decay rate.

The second floor keeps \(D\) fixed but removes proportionality. Now
\(Q^T\Dt Q\) is not diagonal in the graph eigenbasis \(Q\). Off-diagonal modal
damping couples the modes. The eigenvalue problem remains a quadratic
eigenvalue problem (QEP), but graph-frequency movement is no longer a damping
certificate \cite{tisseur2001qep}.

The third floor appears when auxiliary states are eliminated. If the full
linearization is partitioned into retained electromechanical states \(x_e\) and
condensed controller states \(x_c\),
\begin{equation}
  A=
  \begin{bmatrix}
    A_{ee} & A_{ec}\\
    A_{ce} & A_{cc}
  \end{bmatrix},
  \label{eq:intro_partition}
\end{equation}
then Schur condensation gives a nonlinear retained-state operator,
\begin{equation}
  S(s)=sI-A_{ee}-A_{ec}(sI-A_{cc})^{-1}A_{ce}.
  \label{eq:intro_schur}
\end{equation}
In second-order retained coordinates the general condensed operator is
\begin{equation}
  T(s)=s^2M+sD_0+\Lt+\Pi(s),
  \qquad
  \Pi(s)=\Pi_K(s)+s\Pi_D(s).
  \label{eq:intro_sigma}
\end{equation}
The rational correction may contain stiffness-like and damping-like channels.
Only when the stiffness-like part is absent or explicitly absorbed does this
reduce to \(T(s)=s^2M+s\Sigma(s)+\Lt\), with
\(\Sigma(s)=D_0+\Pi_D(s)\). The reduced damping object is then rational in
\(s\), with poles inherited from the eliminated controller block.

The same operator also defines the thesis' action calculus. On the resolvent
set of \(T\), let
\begin{equation}
  G_T(s)=T(s)^{-1}.
  \label{eq:intro_green_function}
\end{equation}
For a line \(e=(i,j)\), with incidence vector \(a_e=e_i-e_j\), define
\begin{equation}
  \mathcal R_e(s)=a_e^TG_T(s)a_e .
  \label{eq:intro_dynamic_resistance}
\end{equation}
At \(s=0\), when the reduced operator has the Laplacian limit and the reference
angle is removed, this recovers classical effective resistance
\(a_e^TL^\dagger a_e\). Away from zero it is a dynamic effective resistance:
it records stiffness, inertia, damping, non-normality, controller poles,
residues, and modal coupling through the same port projection. A finite line
change \(\Delta b_e\) then satisfies
\begin{equation}
  \det(T(s)+\Delta b_ea_ea_e^T)
  =
  \det T(s)\left(1+\Delta b_e\mathcal R_e(s)\right).
  \label{eq:intro_line_action_law}
\end{equation}
Thus a local line sensitivity is only the small-\(\Delta b_e\) limit of an
exact scalar finite-action law.

\begin{figure}[ht]
\centering
\begin{tikzpicture}[
  box/.style={rectangle, rounded corners=2pt, draw=black!65, align=center,
    minimum height=1.06cm, text width=0.72\textwidth},
  level0/.style={box, fill=softblue, draw=thesisblue},
  level1/.style={box, fill=softgreen, draw=thesisgreen},
  level2/.style={box, fill=softorange, draw=thesisorange},
  transition/.style={font=\small, align=center, fill=white, inner sep=2pt},
  arrow/.style={-{Latex[length=2.4mm]}, thick}
]
\node[level0] (l0) {Level 0: commuting fixed damping\\
\([\Lt,\Dt]=0\), including \(D=\gamma M\) and Rayleigh damping};
\node[level1, below=1.15cm of l0] (l1) {Level 1: fixed non-commuting damping\\
\(D\) is constant, but \(Q^T\Dt Q\) couples graph modes};
\node[level2, below=1.15cm of l1] (l2) {Level 2: rational self-energy and Green function\\
\(\Pi(s)=\Pi_K(s)+s\Pi_D(s)\), \(G_T(s)=T(s)^{-1}\)};
\draw[arrow] (l0) -- node[transition,right=0.25cm]{non-commuting damping operator} (l1);
\draw[arrow] (l1) -- node[transition,right=0.25cm]{controller-state condensation} (l2);
\end{tikzpicture}
\caption{Damping-operator hierarchy developed in this thesis. Each transition
keeps part of the previous intuition but removes one mathematical guarantee.}
\label{fig:intro_operator_hierarchy}
\end{figure}

\section{Core Calibrated Messages}

The hierarchy leads to five messages that are strong enough to guide the
thesis, but only if their assumptions are stated precisely.

\begin{enumerate}
\item \textbf{Pure inertia is not damping.} In the fixed-mode, fixed-\(D\)
screen,
\begin{equation}
  \frac{\partial\zeta_k}{\partial M_i}
  \approx
  -\frac{1}{2}\zeta_k\phi_k(i)^2.
  \label{eq:intro_inertia_message}
\end{equation}
Thus adding inertia at a modal antinode can reduce the scalar damping-ratio
screen if damping is not co-designed. This does not say inertia is bad. It
says inertia, damping ratio, RoCoF, and nadir are distinct objectives.

\item \textbf{Graph stiffness is not a damping certificate.} For a line
\(e=(i,j)\),
\begin{equation}
  \frac{\partial\zeta_k}{\partial b_e}
  =
  -\frac{\zeta_k}{2\nu_k}(\phi_k(i)-\phi_k(j))^2
  +
  \frac{\zeta_k}{\delta_k}
  \frac{\partial\delta_k}{\partial b_e}.
  \label{eq:intro_line_message}
\end{equation}
If the line action raises stiffness but does not raise modal damping, the
scalar damping-ratio screen can decrease. Therefore algebraic connectivity and
effective-resistance improvements are graph screens, not automatic
small-signal damping certificates.

\item \textbf{Modal reduction and controller condensation share a resolvent
structure.} Eliminating coupled graph modes from a fixed-\(D\) QEP produces a
rational correction such as \eqref{eq:general_modal_correction}. Eliminating
controller states produces the Schur self-energy
\(\Pi(s)=G(sI-A_{cc})^{-1}(B_\theta+sB_\omega)\). The common mechanism is not
the physical origin of the eliminated states; it is the resolvent injected into
the retained operator. The damping-channel notation \(\Sigma(s)\) is used only
after the rational correction has been projected onto the velocity channel.

\item \textbf{Static damping bridges have a spectral validity region.} In the
damping-channel specialization, the NEP sensitivity denominator contains
\begin{equation}
  T_s(s)=2sM+\Sigma(s)+s\Sigma'(s).
  \label{eq:intro_spectral_memory}
\end{equation}
The term \(s\Sigma'(s)\) is the spectral memory of the condensed controller
block. Setting it to zero is acceptable only when the rational self-energy is
locally flat enough over the critical region. Near a controller resonance, a
static bridge can predict the wrong pole-motion direction.
\item \textbf{Finite actions are not only derivatives.} For line, damping,
inertia, and controller actions that enter \(T(s)\) through low-rank port
updates, the pole condition can be reduced to a determinant whose dimension is
the number of action ports. This does not replace full-model validation, but it
does give the right intermediate object between a scalar graph score and a
large NEP recomputation.
\end{enumerate}

None of these messages is presented as an industry-wide theorem without
qualification. They are calibrated statements: exact in the commuting
constant-\(D\) floor, first-order screens in fixed non-commuting damping, and
validation gates after controller condensation.

\section{Research Questions}

The thesis is organized around eight research questions.

\begin{enumerate}
\item What exact assumptions make graph-mode damping analysis valid in the
classical swing model?
\item How does the operating point induce the operational graph
\(L=B W(z^\star)B^T\), and how do power-flow geometry, line loading, cut
structure, and cycle structure shape dynamic modes?
\item How does fixed non-commuting damping create modal coupling, and how
can that coupling be expressed as a rational correction?
\item How does Schur condensation of controller states transform a fixed
damping matrix into a rational self-energy \(\Pi(s)\), and when is its
damping-channel projection \(\Sigma(s)\) legitimate?
\item How does the dynamic Green function \(G_T(s)\) turn effective
resistance, line reinforcement, damping placement, inertia placement, and
controller retuning into one low-rank action calculus?
\item What theoretical requirements must a planning metric satisfy before it
can claim to identify weak buses, weak lines, or useful converter actions in a
converter-rich grid?
\item How can dynamic hosting capacity, virtual-inertia placement, and
line-reinforcement screening be expressed as spectral margin problems rather
than as static graph or voltage heuristics?
\item Which certificates decide when a reduced QEP, a Schur NEP, a
data-identified self-energy, or a full DAE is adequate for a protected
decision?
\end{enumerate}

The first five questions are mathematical. The last three are epistemic,
planning-oriented, and computational: they define what future empirical
validation must show before a screen becomes an engineering certificate.

\section{Contributions}

The dissertation makes one central contribution and five supporting
contributions.

\begin{enumerate}
\item \textbf{Central contribution: operator-valued dynamic grid strength.}
The dissertation treats graph stiffness, damping-to-inertia roughness,
spectral effective resistance, rational controller self-energy, and nonlinear
memory as successive mechanisms in one retained operator
\(T(s)=s^2M+sD+L+\Pi(s)\).
\item \textbf{Damping-to-inertia graph roughness.} It shows how
non-commutation between inertia-normalized stiffness and damping becomes the
graph-gradient energy of \(D_i/M_i\), and how this roughness, together with
small modal gaps, predicts modal mixing.
\item \textbf{Graph-time resonance atlas.} It defines
\(H_G(\omega)=[Q^T\widetilde T(j\omega)Q]^{-1}\) as the joint spatial-mode
and temporal-frequency response used to locate resonant modal conversion.
\item \textbf{Spectral effective resistance.} It defines
\(\mathcal R_e(s)=a_e^TG_T(s)a_e\) and derives the finite line-update law
\(1+\Delta w_e\mathcal R_e(s)=0\), recovering static effective resistance at
the Laplacian limit.
\item \textbf{Controller-induced reinforcement reversal.} It formalizes the
case in which a line or parameter action passes a static graph or fixed-\(D\)
screen but reverses once the controller Schur resolvent is included.
\item \textbf{Certificate/screen/gate planning methodology.} It formulates
weak elements, dynamic hosting, virtual-inertia placement, and line
reinforcement as operator-, margin-, action-, and validation-gate-relative
claims rather than graph-intrinsic labels.
\end{enumerate}

\section{Dissertation Organization}

The dissertation is now organized as a modular monograph. Part~I builds the
problem statement, the physical grid framework, the spectral preliminaries, and
the state of the art. The new physical framework chapter develops the AC
network, phasor, power-flow, swing-equation, DAE, and operator chain before the
more compact mathematical chapters use those objects. Part~II develops the
operational graph theory: power-flow geometry, total derivatives of
\(L=B W(z^\star)B^T\), graph roughness of damping-to-inertia fields,
graph-time response, dynamic effective resistance, and structured fragility
radii. Part~III closes the constant-\(D\) theory: QEP identities, commuting
damping, uniform decay, Rayleigh damping, non-commuting coupling, and
reduced-QEP gates. Part~IV develops rational self-energy \(\Pi(s)\), its
damping-channel specialization \(\Sigma(s)\), Schur condensation,
finite-action laws, NEP sensitivities, port equivalence, passivity screens,
structured fragility, and controller-resolvent reversals. Part~V separates
genuine nonlinear DAE stability from spectral NEP theory by treating
equilibria, nonlinear memory, bifurcation, energy, and regions of attraction.
Part~VI turns the operators into applications: weak elements, planning risk
vectors, dynamic hosting regions, reinforcement envelopes, and reproducible
benchmark protocols.

The integrated source-to-claim ledger is not part of the main scientific
progression. It is kept as an appendix-level traceability record that maps
claims to sources, scripts, figures, and validation gates.
